\bookchapter{פרק 18: גיבוי והעברת נתונים}

\chapterintro{לנתונים שונים יש “מצבי הישרדות” שונים, צריך להבדיל ביניהם:}

\booksection{נשמרים מקומית בטלפון, אובדן אמיתי}

: תמונות וסרטונים (אם לא הופעל סנכרון ענן), פתקים, SMS, היסטוריית שיחות, קבצים שהורדו.

\begin{itemize}
  \item \textbf{נשמרים בשרתי אפליקציה, החלפת טלפון והתחברות מחדש מחזירה}: רשימת אנשי קשר ב-WeChat, הזמנות ב-Taobao, היסטוריית עסקאות באפליקציית בנק. הדברים האלה מאוחסנים בשרתים של החברה המתאימה, לא בטלפון שלך.
  \item \textbf{אזור אפור, דורש פעולה מיוחדת}: היסטוריית צ'אט WeChat (למרות ששרת WeChat שומר, אי אפשר להוריד מהשרת בעצמך), אישורי מכשיר של בנק מקוון, מפתחות באפליקציות אימות דו-שלבי (כמו Google Authenticator).
  \item \textbf{לכן מה שצריך להגן עליו במיוחד הוא הסוג הראשון והשלישי.}שתיים, גיבוי מלא: הכל במכה אחת
\end{itemize}

גיבוי מלא הוא כמו לצלם “תמונת מצב” מלאה של הטלפון - הגדרות מערכת נוכחיות, פריסת שולחן העבודה, נתוני אפליקציות, תמונות, אנשי קשר - הכל נארז.

\booksection{תוכנית א': גיבוי מקומי במחשב (האמין ביותר, מומלץ למי שיש מחשב)}

iPhone

\booksubsection{: חבר את הטלפון בכבל למחשב, פתח את iTunes (גרסת Windows ישנה) או Finder (Mac), בדף המכשיר בחר “גבה עכשיו”. שים לב לסמן “הצפן גיבוי”, אחרת סיסמאות ונתוני בריאות לא יגובו.}

\begin{itemize}
  \item \textbf{אנדרואיד}: לכל מותג יש תוכנת מחשב משלו, כמו Huawei Phone Assistant, Xiaomi Assistant, Samsung Smart Switch. חבר למחשב ופעל לפי ההוראות.
  \item \textbf{HarmonyOS}תוכנית ב': גיבוי ענן (הכי נוח, מומלץ לכולם)
  \item iPhone
\end{itemize}

\booksubsection{: הגדרות השם שלך iCloud גיבוי ענן iCloud הפעל. הטלפון יגבה אוטומטית כשמחובר לחשמל, ל-Wi-Fi, ונעול. 5GB החינמיים כנראה לא מספיקים, מומלץ לשלם 6 יואן ולשדרג ל-50GB.}

\begin{itemize}
  \item \textbf{אנדרואיד}: לכל מותג יש פונקציית גיבוי ענן. Huawei בהגדרות חפש "גיבוי ענן", Xiaomi ב-"שירות ענן Xiaomi", OPPO/vivo דומה. בנוסף יש גיבוי Google מובנה (הגדרות מערכת גיבוי), גם אותו כדאי להפעיל.
  \item \textbf{HarmonyOS}שלוש, אסטרטגיית גיבוי מסווג: דברים חשובים מגובים בנפרד
  \item גיבוי מלא הוא רשת ביטחון, אבל יש דברים חשובים מדי, שווים עותק נפרד נוסף:
\end{itemize}

\booksection{תמונות וסרטונים (הנכס הדיגיטלי היקר ביותר)}

תוכנית ביטוח כפול

\booksubsection{: הפעל סנכרון אוטומטי בתמונות iCloud/Google תמונות, ובנוסף מדי כמה חודשים העתק ידנית תמונות למחשב או כונן קשיח נייד. כך עותק בענן, עותק מקומי, משמשים כגיבוי הדדי.}

\begin{itemize}
  \item \textbf{תוכנית בלי תשלום}: Google תמונות נותנת 15GB בחינם, אפשר להגדיר "איכות דחוסה" לחיסכון במקום. או מדי חודש להעתיק את כל התמונות בכבל למחשב.
  \item \textbf{אנשי קשר (קל להתעלם, אבל חיוני ביותר)}אנשי קשר ב-iPhone מסונכרנים כברירת מחדל עם iCloud, אם הופעל סנכרון אין צורך בגיבוי נוסף. אבל ליתר ביטחון, אפשר לייצא קובץ vCard מ-iCloud.com למחשב.
\end{itemize}

\booksubsection{אנשי קשר באנדרואיד בדרך כלל מסונכרנים עם Google או ענן המותג. בהגדרות אפליקציית אנשי הקשר ייצא קובץ .vcf, זהו גיבוי מלא.}

\begin{itemize}
  \item היסטוריית צ'אט WeChat (טיפול נפרד, החשוב מכל)
  \item חובה לדבר על זה בנפרד, כי בגיבוי מלא היסטוריית הצ'אט של WeChat לרוב לא מלאה, וכלי העברת מכשיר בכלל לא מצליחים להזיז אותה.
\end{itemize}

\booksubsection{שיטה א': העברה לטלפון חדש}

. WeChat אני הגדרות כללי העברה וגיבוי היסטוריית צ'אט העברה העברה לטלפון/טאבלט. הטלפון הישן והחדש מתחברים לאותו Wi-Fi, סורקים קוד ומעבירים. שים לב: לפני סיום ההעברה, אל תצא מ-WeChat בטלפון הישן.

\begin{itemize}
  \item \textbf{שיטה ב': גיבוי למחשב}. התקן את WeChat PC במחשב, שלושה קווים בפינה השמאלית התחתונה העברה וגיבוי גיבוי ושחזור גיבוי היסטוריית צ'אט למחשב. אחר כך בהחלפת טלפון אפשר לשחזר לטלפון החדש.
  \item \textbf{שיטה ג': גיבוי לענן}. WeChat תומך גם בגיבוי ל-iCloud או לענן הטלפון, אבל שים לב: הגיבוי הזה ניתן לשחזור רק בטלפון הנוכחי, לא חוצה מכשירים. בהחלפת טלפון עדיין צריך להשתמש בשתי השיטות הקודמות.
  \item \textbf{ארבע, מדריך החלפת מכשיר: מהטלפון הישן לחדש}החלפת טלפון חדש זה שמח, אבל אם הנתונים לא עוברים זה הופך לכאב ראש. שלושה מקרים, כל אחד בנפרד:
\end{itemize}

\booksection{שדרוג באותו מותג (הכי פשוט)}

iPhone iPhone חדש

\booksubsection{: בהפעלה שים את הטלפון הישן ליד החדש, יופיע אוטומטית “התחלה מהירה”, סרוק עם הישן את התבנית על החדש. כל הנתונים, ההגדרות, פריסת האפליקציות יעברו כמו שהם, כולל סיסמאות Wi-Fi. לוקח כ-20-30 דקות.}

\begin{itemize}
  \item \textbf{אנדרואיד באותו מותג}: השתמש בכלי העברת המכשיר של המותג. Huawei משתמש ב"שיבוט טלפון", Xiaomi ב"החלפת מכשיר Xiaomi", OPPO ב"מעבר טלפון", vivo ב"העברה הדדית". השלבים דומים: התקן את אותה אפליקציה בשני הטלפונים, סרוק וחבר, סמן מה להעביר, התחל.
  \item \textbf{החלפת מותג (אנדרואיד לאנדרואיד, מותגים שונים)}השתמש בכלי העברה של צד שלישי או באפליקציית ההעברה של המותג הישן (רוב האפליקציות תומכות בהעברה לאנדרואיד אחר). מה עובר: אנשי קשר, תמונות, SMS, היסטוריית שיחות, חלק מהאפליקציות. מה לא עובר: מצב התחברות בתוך אפליקציות (צריך להתחבר מחדש), היסטוריית צ'אט WeChat (השתמש בפונקציית ההעברה של WeChat עצמה).
\end{itemize}

\booksubsection{החלפת פלטפורמה (אנדרואיד ל-iPhone, או להפך)}

אנדרואיד iPhone

\booksubsection{: ל-Apple יש אפליקציית “Move to iOS”, מורידים באנדרואיד, בהפעלת ה-iPhone החדש בוחרים “העבר נתונים ממכשיר אנדרואיד”. עוברים אנשי קשר, תמונות, SMS ונתונים בסיסיים, אבל נתוני אפליקציה לא. היסטוריית צ'אט WeChat עדיין דורשת את פונקציית ההעברה של WeChat.}

\begin{itemize}
  \item \textbf{iPhone אנדרואיד}: לכל מותג אנדרואיד יש כלי העברה משלו, אבל להגיד את האמת, חוויית ההעברה חוצת פלטפורמות לא טובה כמו באותה פלטפורמה. הדברים החשובים ביותר קודם: ייצוא אנשי קשר, גיבוי תמונות, העברת היסטוריית צ'אט WeChat. השאר לאט לאט.
  \item \textbf{חמש, חוק הזהב של גיבוי: עקרון 3-2-1}זהו עיקרון קלאסי בתחום אבטחת הנתונים, פשוט לזכירה:
\end{itemize}

\booksection{3}

: שמור לפחות שלושה עותקים של נתונים חשובים (מקור אחד, שני גיבויים).

\begin{itemize}
  \item \textbf{2}: השתמש בשני אמצעי אחסון שונים לגיבוי (למשל עותק בענן, עותק בכונן קשיח נייד).
  \item \textbf{1}: לפחות עותק גיבוי אחד במיקום פיזי אחר (למנוע מצב שאש, גניבה וכו' יהרסו את כל הנתונים בבת אחת).
  \item \textbf{לאדם מן השור, "עותק בענן, עותק במחשב" כבר נחשב מצוין. לא צריך לסבך, העיקר שבאמת תעשה.}שים לב מיוחד להיסטוריית צ'אט WeChat
\end{itemize}

לא ניתן להעביר היסטוריית צ'אט WeChat עם אף כלי העברת מכשיר! היא לא בטווח הכיסוי של גיבוי מלא. חובה להשתמש בפונקציית “העברת היסטוריית צ'אט” המובנית של WeChat, או לגבות ל-WeChat למחשב. הרבה אנשים מגלים אחרי החלפת הטלפון שכל היסטוריית הצ'אט נעלמה, ואז כבר מאוחר - הטלפון הישן אולי כבר עבר איפוס להגדרות יצרן.

\begin{tcbwarn}{תתאמן: תרגול גיבוי מלא}
בדוק שגיבוי ענן iCloud או גיבוי ענן המותג מופעל, אשר את זמן הגיבוי האחרון.

\end{tcbwarn}

\begin{tcbpractice}{גבה תמונות ידנית: חבר בכבל למחשב, העתק את תיקיות DCIM ו-Pictures למחשב.}
\begin{enumerate}
  \item גבה אנשי קשר: משתמש iPhone ייצא vCard מ-iCloud.com; משתמש אנדרואיד ייצא קובץ .vcf מהגדרות אפליקציית אנשי הקשר.
  \item גבה היסטוריית צ'אט WeChat: פתח WeChat אני הגדרות כללי העברה וגיבוי היסטוריית צ'אט גיבוי למחשב.
  \item העתק את קובץ ה-vcf של אנשי הקשר ואת תיקיית התמונות לכונן קשיח נייד או USB, להשלמת "עותק אחד לא מקוון במקום אחר" שבכלל 3-2-1.
  \item פרק 19: אבחון תקלות והצלה עצמית
  \item הטלפון, כמו בני אדם, מדי פעם "חולה". קפיאות, קריסות, מסך שחור, לא טוען... רוב התקלות לא באמת נוראיות, אפשר לנסות קודם בעצמך לפני שליחה לתיקון. הפרק הזה הוא "מדריך מיון" לטלפון שלך - קודם נסה שיטות זולות, ורק אם לא עוזר, פנה לרופא מקצועי.
\end{enumerate}

\end{tcbpractice}
